\scanpage{35}
\mspage{34}

% slanted
\margin{axiome trivial}
\begingroup\makeatletter\tagsleft@true\makeatother
\[
 X\leq Y\quad\Longrightarrow\quad X\ll Y
 \tag*{$(\deleted{\unsure{C|G}}\mathrm{ML\,0})$}
\]

% slanted
\margin{axiome du raffinement induit}
$(\deleted{C}\mathrm{ML\,1})$ \deleted{\unsure{ent}}\unsure{Soit|Soient} $\deleted{\unsure{X'}}$
\diagram{\begin{tikzcd}[ampersand replacement=\&,row sep=0pt,column sep=3pt]
 Z \& \ll \& X\\
 \& \& \rotatebox{90}{$\leq$}\\
 \& \& X'
\end{tikzcd}}
dans $M$, et soit $x\in\unsure{L|\mathcal L}$ avec $x\ll Z$, $x\ll X'$. Alors $\exists\,Z'\in M$ avec
\diagram{\begin{tikzcd}[ampersand replacement=\&,row sep=1pt,column sep=4pt]
 \& \& Z \& \ll \& \unsure{x|X}\\
 \& \& \rotatebox{90}{$\leq$} \& \& \rotatebox{90}{$\leq$}\\
 x \& \ll \& Z' \& \ll \& X'\rlap{$\quad\unsure{,|.}$}
\end{tikzcd}}

\uline{De} plus, si $Z\leq X$, on peut supposer $Z'\leq X'$.

% slanted
\margin{axiome de divisibilité \deleted{\unsure{pour|par}}}
\[
\begin{aligned}
 &\forall X\in M,\quad \unsure{\mathrm{omb}|\mathrm{ends}|\mathrm{ens}}(X)^{\unsure{\circ|0}}\ne 0
   \quad\text{i.e. }\exists\,x\in L,\ \text{avec}\\
 &x\mathrel{\overset{\circ}{\ll}}X\quad\text{i.e. }x\ll X\text{ et }
   (\unsure{a|x|\alpha}\ll X'\leq X_{\unsure{x|n}})\Longrightarrow X'=X
\end{aligned}
\tag*{$(\mathrm{ML\,2})^{*}$}
\]

% slanted
\margin{axiome de \deleted{\unsure{plus}} la \deleted{plus \unsure{petite}} strate minimale}
\[
\begin{aligned}
 &\deleted{\unsure{\forall x\in}}\ \forall\,X\in M,\ x\unsure{\in|\notin}L
    \ \text{avec }\ x\ll X,\ \exists\,X'\leq X\\
 &\text{avec }\quad x\mathrel{\overset{\circ}{\ll}}X'\,.
\end{aligned}
\tag*{$(\mathrm{ML\,3})$}
\]

$(\mathrm{ML\,4})$ $\parallel$ est \unsure{sym,|sym.} antiréflexive \inserted{$(x\parallel y\Rightarrow y\parallel x\text{ et }x\ne y)$}

% slanted
\margin{C'est le seul des axiomes qui \unsure{fasse} intervenir la relation $\parallel$}
\margin{$\swarrow\!\big\lfloor$ axiome de disjonction pour les lieux}
\[
\left\{\begin{aligned}
 &(\mathrm{ML}\,\deleted{\unsure{4}}5)\qquad \forall\,X,Y\in M\ \text{avec }\ Y\leq X,\quad\deleted{\text{\unsure{si }}y}\\
 &\hspace{3.5em} y\in\mathrm{omb}(Y),\ x\in\unsure{\mathrm{om}|\mathrm{en}|\mathrm{omb}}(X)
    \setminus\mathrm{omb}(Y)\quad\Longrightarrow\quad x\parallel y
\end{aligned}\right.
\]
\endgroup

\deleted{Les atlas correspondant à une telle donnée, réalisés comme des ens de parties de $M$}
\[
\deleted{\unsure{\mathcal F|F}\subset\unsure{\mathrm P|\mathcal P}(M),}
\]
\deleted{sont \deleted{\unsure{toute|toutes} \unsure{forme}} les $\mathcal F$ satisfaisant les conditions suivantes}

\deleted{(i)}

$*$ \uline{NB} Cet axiome implique que si \deleted{\unsure{la}}les relations d'ordre $\leq$ et $\ll$ sont égales, alors $\deleted{\unsure{X}}M$ est discret i.e. $M=L$. \deleted{\unsure{La thèse est vide}} L'axiome $\mathrm{ML\,4}$ est vide !
